\section{Results and Discussion}
\label{sec:results}
\rev{The operating point} is affected by a good choice of the cutoff threshold on $P(\mathrm{malignant})$. A case with $P(\mathrm{malignant})$ $\geq$ cutoff is considered malignant otherwise it is considered benign. The threshold is chosen on the validation images only\rev{, from 0.05 to 0.95 in steps of 0.01, and is kept unchanged for the test images. Changing the threshold does not change the trained model; it only moves its operating point between sensitivity and specificity.} Table~\ref{tab:comparison} shows the effect of \rev{the threshold rule} on the test performance. \rev{A designer should select the operating point that matches the clinical goal.} For breast cancer diagnosis one should target the sensitivity metric since it avoids missing a positive case.



Table~\ref{tab:comparison} compares the selected model (first row in Table~\ref{tab:comparison}) with the base model on the 644 test abnormalities. The ROC-AUC values are $0.752$ for DeepProbLog and $0.746$ for the baseline. \rev{The 95\% bootstrap intervals are $[0.712, 0.789]$ and $[0.707, 0.784]$. The paired DeLong test gives a difference of $0.005$, with a 95\% interval of $[-0.035, 0.046]$ and $p=0.79$, so the two models rank the test cases equally well. These results come from a single training run.} \rev{\todo{add the extra-seed results when the seed 7 run finishes}} \rev{On validation, the recorded class receives a probability of at least 0.5 for} 90.5\% \rev{of the images for} the lesion kind head, 72.3\% \rev{for} the shape head, 70.1\% \rev{for} the margin head, and 82.0\% \rev{for} the calcification head. Sensitivity and specificity differ because each row uses a different threshold. The lower specificity of the selected model is the cost of the higher sensitivity. \rev{Within the swept range, no baseline threshold reaches a validation sensitivity of 0.80. A lower threshold would reach it: on the test ROC curve, the baseline needs a threshold of 0.022 to reach a sensitivity of 0.80, where its specificity is 0.540, against 0.556 for DeepProbLog at the same sensitivity. At matched sensitivity the two models are therefore close.}

\begin{table}[!t]
\caption{Test comparison on 644 images.}
\label{tab:comparison}
\centering
\footnotesize
\setlength{\tabcolsep}{4pt}
\begin{tabular}{@{}lcccccc@{}}
\toprule
Model & Valid Thr. & Acc. & Sens. & Spec. & Macro-F1 & AUC \\
\midrule
DeepProbLog & 0.55$^{*}$ & 0.654 & \textbf{0.798} & 0.558 & 0.654 & \textbf{0.752} \\
DeepProbLog & 0.60$^{\dagger}$ & 0.666 & 0.716 & 0.633 & 0.663 & \textbf{0.752} \\
ResNet18 & 0.50 & 0.689 & 0.591 & \textbf{0.755} & 0.674 & 0.746 \\
ResNet18 & 0.44$^{\ddagger}$ & \textbf{0.693} & 0.611 & 0.747 & \textbf{0.679} & 0.746 \\
\bottomrule
\end{tabular}

\vspace{0.45em}
\begin{minipage}{\columnwidth}
\footnotesize
$^{*}$\rev{Among validation thresholds with sensitivity at least 0.80, the one with the highest validation specificity.}\\
$^{\dagger}$Validation threshold with the highest macro-F1.\\
$^{\ddagger}$The baseline never reached sensitivity 0.80 on validation \rev{within the swept range}. This is the threshold that best balances its sensitivity and specificity.\\
\rev{AUC 95\% bootstrap CI: DeepProbLog $[0.712, 0.789]$, ResNet18 $[0.707, 0.784]$; DeLong $p=0.79$.}\\
Valid Thr.: validation threshold. Acc.: accuracy. Sens.: sensitivity. Spec.: specificity. AUC: area under the curve, which does not depend on a fixed threshold.
\end{minipage}
\end{table}



\subsection{\rev{Why the Gain Is Small}}
\label{sec:whysmall}

A gain of $0.006$ in ROC-AUC is not a large improvement\rev{, and the DeLong test shows it is not significant. At both thresholds the accuracy and macro-F1 of DeepProbLog are lower than those of the baseline.} Two limitations are behind this result:
\begin{itemize}
\item Both systems start from the same ResNet18  architecture pretrained on ImageNet. Next, they are  finetuned on the same $224\times224$ crops.

\item \rev{The three rules do not cover every malignant case. For example, patient P\_00017 has a malignant mass with a round shape, recorded as regular. Its margin is circumscribed in one view and ill-defined (indistinct) in the other, which falls into other. No rule matches either view even with correct concepts. This is a gap in the rules, not label noise: an indistinct margin is suspicious in BI-RADS.}

\end{itemize}

\begin{revblock}
\textit{Ceiling of the rules.} To separate the rules from the heads, the program was applied to the recorded concepts of the 644 test images, without the network (Table~\ref{tab:ceiling}). With these certain facts, the three rules reach a sensitivity of 0.805 and a specificity of 0.587. Extending the margin rules to indistinct and microlobulated margins, and treating a missing calcification type as unknown, raises the sensitivity to 0.918 but lowers the specificity to 0.468. ROC-AUC moves from 0.696 to 0.693, so the trained rules were left unchanged: the comparison is on ROC-AUC, and the extension does not improve it. Indistinct margins are suspicious in BI-RADS and common (24.5\% of training mass crops, 30.6\% of test mass crops); they stay in other because the added rule does not improve ranking. Microlobulated margins are uncommon (8.5\% and 6.4\%), so a separate trained class would be small, and the same check shows no ROC-AUC gain. The margin head was not retrained on these classes. That retraining, and calcification distribution, are left for future work.

\textit{Rules or heads.} At the 0.55 threshold, DeepProbLog makes 223 errors on the test images. In 130 of them (58\%; 21 missed malignancies and 109 false alarms), the rules give the wrong answer even on the recorded concepts, so the error comes from the rules. In the other 93 (42\%; 31 missed malignancies and 62 false alarms), the rules would give the correct answer on the recorded concepts, so the error comes from the heads.

\todo{Stage 1 alone vs.\ Stage 2 on the paper weights, and test accuracy and confusion matrices of the four heads (evaluation of the saved checkpoints on Kaggle)}
\end{revblock}

\begin{table}[!t]
\color{revcolor}
\caption{The three rules applied to the recorded concepts of the 644 test images (no network).}
\label{tab:ceiling}
\centering
\footnotesize
\setlength{\tabcolsep}{4pt}
\begin{tabular}{@{}lccccc@{}}
\toprule
Margin rules & Acc. & Sens. & Spec. & Macro-F1 & AUC \\
\midrule
Spiculated (as trained) & 0.674 & 0.805 & 0.587 & 0.674 & 0.696 \\
+ indistinct, microlobulated & 0.648 & 0.918 & 0.468 & 0.645 & 0.693 \\
\bottomrule
\end{tabular}
\end{table}


\subsection{A Decision Trace}
\label{sec:traces}

We consider the case of a benign mass (P\_00061) (whose ground truth description is: lesion kind= mass, shape= regular, margin = circumscribed, and calcification= absent) that has been misclassified by the model to highlight the practical utility of the decision trace. A standard black-box network would merely output a highly misleading malignancy score, leaving us with no clue to detect the error. Table~\ref{tab:trace}, gives the logic program findings with their corresponding probabilities. The root cause of this failure is that the neural head "lesion kind" misidentified the mass as a calcification (with a probability of 0.99\rev{),} which is multiplied by $P(\mathrm{calcification}=\mathrm{suspicious})=0.919$ \rev{in Eq.~(\ref{eq:query}) to yield $P(\mathrm{malignant})=0.91$. The mass branch adds almost nothing, because $P(\mathrm{mass})=0.010$.} This allows us to spot the contradictory evidence, reject the erroneous conclusion, and understand exactly where the model's logic failed. This information is, also, valuable for the trainer/ designer of the architecture to improve the network head causing the error.

\begin{table}[!t]
\caption{Trace for benign mass P\_00061. \rev{Only the classes used by the rules are shown.}}
\label{tab:trace}
\centering
\footnotesize
\setlength{\tabcolsep}{4pt}
\begin{tabular}{@{}lr@{}}
\toprule
Predicted Class & Value \\
\midrule
$P$(kind $=$ mass) & 0.010 \\
$P$(kind $=$ calcification) & 0.990 \\
\rev{$P$(shape $=$ irregular)} & \rev{0.001} \\
\rev{$P$(margin $=$ spiculated)} & \rev{0.000} \\
$P$(calcification $=$ suspicious) & 0.919 \\
$P$(calcification $=$ absent or typically benign) & 0.081 \\
\midrule
$P$(malignant) & 0.910 \\
\bottomrule
\end{tabular}
\end{table}
